\scanpage{6}
\mspage{3}
\[
(14)\qquad
1\mathrel{\vcenter{\hbox{\begin{tikzpicture}[x=1mm,y=1mm]
\draw[line width=.45pt] (0,0)--(3,0);
\draw[line width=.6pt] (3,.1)--(3,-.22);
\end{tikzpicture}}}}\pi_1(X')
\mathrel{\vcenter{\hbox{\begin{tikzpicture}[x=1mm,y=1mm]
\draw[line width=.45pt] (0,0)--(6,0);
\draw[line width=.65pt] (6,.24)--(6,-.24);
\end{tikzpicture}}}}\mathcal E_{X'}
\longrightarrow\mathcal G_{X'}\longrightarrow1.
\]
et on a encore un germe de \deleted{hom} \inserted{relèvement}
\unsure{Circonstanciellement|circonstanciellement} mod comp. par autom int.
provenant de $\pi_1(X')$)
\diagram{$(15)\qquad$
\begin{tikzcd}[ampersand replacement=\&]
\mathcal E_{X'}\arrow[r,phantom,
"{\vcenter{\hbox{\begin{tikzpicture}[x=1mm,y=1mm]
\draw[line width=.45pt] (0,0)--(6,0)--(6,.5);
\end{tikzpicture}}}}" description]
\&\mathcal G_{X'}\arrow[r]\&\Gamma\arrow[d,phantom,"\cup" description]\\
\&\&\Gamma'\arrow[ull,"{\varphi_\xi}"{below}]
\end{tikzcd}}
(si $\xi\in X'(\overline{\mathbf Q})$).

Notons qu'on a ici un hom
\[
(16)\qquad\mathcal G_{X'}\longrightarrow\mathfrak S_I
\]
La théorie de Kummer nous fournit,
pour $\forall\,i$, un hom extérieur
\diagram{%
\begin{tikzpicture}[x=1cm,y=1cm,line width=.45pt]
\node at (-2.45,-.64) {$(17)$};
\node (t17) at (0,0) {$\mathbf T(\overline{\mathbf Q})$};
\node at (4.75,0) {$\pi_1(X')$};
\draw[->] (1.08,0)--(3.83,0);
% The phi_i label's own long stem crosses the horizontal shaft.
\draw (2.45,-.25)--(2.45,.61)
 .. controls (2.29,.73) and (2.43,.97) .. (2.56,.99);
\draw (2.45,.61)
 .. controls (2.56,.55) and (2.72,.65) .. (2.68,.80)
 .. controls (2.63,.94) and (2.46,.86) .. (2.45,.61);
\node[anchor=west] at (2.75,.53) {$\scriptstyle i$};
\draw[->] (-.10,-.34)
 .. controls (-.26,-.58) and (-.18,-.89) .. (.16,-.87);
\node[anchor=base west] at (-.44,-1.22) {$=$};
\node[anchor=base west] at (.02,-1.22)
 {$\displaystyle\varprojlim\mu_n(\overline{\mathbf Q})$};
\end{tikzpicture}}
\margin{$\varphi_i\in\operatorname{Homext}
(\mathbf T(\unsure{\mathbf Q|\overline{\mathbf Q}}),\pi_1(X'))$}
{\reversemarginpar
\margin{\deleted{il faut supposer $\chi=1$
\unsure{si l'on veut}
\unsure{interpréter ça par} \illegible{?}
\unsure{purement arithmétique}
\unsure{un objet} \unsure{\dots}}}
\normalmarginpar}

Notons aussi un hom canonique
\[
(18)\qquad
\left(\mathcal G_{X'}
\mathrel{\vcenter{\hbox{\begin{tikzpicture}[x=1mm,y=1mm]
\draw[line width=.45pt] (0,0)--(5.5,0)--(5.55,.22)--(5.7,0);
\end{tikzpicture}}}}\right)
\Gamma\xrightarrow{\chi}\widehat{\mathbb Z}^{*}
\xrightarrow{\sim}\operatorname{Aut}(\mathbf T(\overline{\mathbf Q}))
\]
faisant opérer $\mathcal G_{X'}$ via $\Gamma$ sur $\mathbf T(\overline{\mathbf Q})$.
Ceci dit, l'\unsure{opération|Opération} extérieure de $\mathcal G_{X'}$
sur $\pi_1(X')$ est compatible avec la famille
\mspage{5}
